\section{The Gaussian sign posterior: negative exponents}
\label{sec:negative-gamma}

Let $G\sim\mathcal N(0,K)$ on the finite cube, and condition on
\[
 \mathcal A_S=\{f(X_i)G(X_i)>0:\ 1\le i\le n\}.
\]
Write $m=\mathbb E[G\mid\mathcal A_S]$ and
$q_S(x)=\mathbb P\{f(x)G(x)>0\mid\mathcal A_S\}$.
We use the buffered threshold
\[
 \theta=\frac{13}{20},\qquad
 c_\theta=\sqrt{2\pi}\left(\frac1e-\frac7{20}\right)>0,
 \qquad F_\theta(S)=\mathbb P_x\{q_S(x)\le\theta\}.
\]
The buffer $\theta<2/3$ will preserve the strict confidence convention
when taking the finite-width limit.

\begin{lemma}[Gaussian mean and confidence]
\label{lem:negative-mean-confidence}
If $s$ is the number of distinct training inputs, then
\[
 \|m\|_K^2\le s\le n,\qquad f(X_i)m(X_i)>0.
\]
Let $\sigma_x^2=\operatorname{Var}(G(x)\mid G_S)$, with $S$ listing
the distinct inputs. If $\sigma_x>0$ and $q_S(x)>\theta$, then
$f(x)m(x)>c_\theta\sigma_x$.
For every Fourier cutoff $v$, moreover,
\begin{equation}\label{eq:negative-variance-tail}
 \mathbb P_x\{\sigma_x^2<\tau_v/2\}
       \le\frac{2n\ell_v}{\tau_v},
 \qquad
 \|m-\Pi_{\le v}m\|_2^2\le n\ell_v.
\end{equation}
\end{lemma}

\begin{proof}
Whiten the training vector as $G_S=K_{SS}^{1/2}\xi$.
The sign constraints form a homogeneous cone in $\xi$. Since a standard
Gaussian's radius is independent of its direction, conditioning leaves
$\mathbb E\|\xi\|_2^2=s$. Jensen's inequality and Gaussian regression
therefore give
\[
 m_S^{\mathsf T}K_{SS}^{-1}m_S\le s,
 \qquad m=K_{\mathcal X,S}K_{SS}^{-1}m_S.
\]
The squared RKHS norm of the regression vector is the first quadratic
form, and its signed training values are positive.

The conditional scalar variable $Y=f(x)G(x)$ has a log-concave density:
it is a marginal of a Gaussian restricted to a convex set, by
Pr\'ekopa's marginal theorem \cite{prekopa1973}. Its cdf $H$ is also
log-concave. For $\mu=\mathbb E Y$, Jensen's inequality and the probability
integral transform give
\[
 \log H(\mu)\ge\mathbb E\log H(Y)=-1,
 \qquad \mathbb P\{Y>\mu\}\le1-e^{-1}.
\]
Conditioning first on $G_S$ expresses the density of $Y$ as a mixture
of Gaussian densities of common variance $\sigma_x^2$. It is bounded
above by $(\sqrt{2\pi}\sigma_x)^{-1}$. If $q_S(x)>\theta$, this first
forces $\mu>0$, and then
\[
 \theta<q_S(x)\le1-e^{-1}
                  +\frac{\mu}{\sqrt{2\pi}\sigma_x}.
\]
This proves the confidence margin.

To bound residual variances, write $G=G_{\le v}+G_{>v}$ independently.
Conditioning additionally on the entire low-degree field can only
reduce a Gaussian conditional variance. The remaining covariance is
\[
 Q_v-Q_{v,\mathcal X,S}Q_{v,SS}^{\dagger}Q_{v,S,\mathcal X}.
\]
The subtracted matrix is positive semidefinite, bounded by $Q_v$, and
has rank at most $s$. Its normalized trace is at most $n\ell_v$.
Markov's inequality proves the first bound in
\eqref{eq:negative-variance-tail}; this automatically includes the
observed points, where the full residual variance vanishes.
The second bound follows from Fourier orthogonality and
$\|m\|_K^2\le n$.
\end{proof}

\begin{theorem}[An exponential lower bound for the sign posterior]
\label{thm:negative-all-n}
For every fixed depth $L\ge2$, there is $c_L>0$ such that, for every
sufficiently large $d$ and each integer $0\le n\le e^{c_Ld}$,
\[
 \mathbb P_S\{F_{13/20}(S)\ge3/8\}\ge3/4.
\]
Consequently the same lower bound holds for the $\liminf$ of the
finite-width strict-$2/3$ confidence-failure fraction, for each fixed
$\gamma<0$. The width limit is taken at fixed $d$ and training sample.
\end{theorem}

\begin{proof}
We give the cutoff estimates explicitly. Fix the sampling-failure
probability at a sufficiently small universal constant and let $A$ be
the resulting constant in Lemma~\ref{lem:parity-squared-balance}.
Choose the largest odd $r$ such that
\[
 n\ge A D_r\log(2D_r).
\]
First suppose $n\ge\exp((\log(d+1))^3)$. Then the adaptive cutoff
estimates of Lemma~\ref{lem:adaptive-cutoff} give
$r\ge(\log d)^2/4$ for large $d$. By reducing $c_L$ we may ensure
$r\le\alpha d$, for a fixed $\alpha<1/256$ to be chosen below.
Write $\rho=r/d$ and $L_\rho=\log(e/\rho)$.
Maximality and the sharp paired eigenvalue estimates give
\begin{equation}\label{eq:negative-cutoff-bounds}
 \frac{n\ell_r}{\tau_r}\le C_A L_\rho,
 \qquad
 \|Q_{r,SS}\|_{\mathrm{op}}\le b_r\tau_r,
 \qquad b_r=C_{L,A}\left(1+\sqrt{L_\rho/\rho}+L_\rho\right),
\end{equation}
with probability at least $1-\delta$, where $\delta$ is a sufficiently
small fixed constant. Indeed,
$n<A D_{r+2}\log(2D_{r+2})$,
$D_{r+2}\asymp\binom{d+1}{r+2}$, and
$\log D_{r+2}\le C r\log(ed/r)$. The ratio estimate
\eqref{eq:ratio} proves the first bound. For the Gram bound use
Lemma~\ref{lem:gram-general}, with
\[
 \frac{nd\,s_{2,r}}{\tau_r^2}\le C_A L_\rho/\rho,
 \qquad \frac{n\ell_r}{\tau_r}\le C_A L_\rho.
\]
The far-pair and collision errors vanish uniformly in this range when
$c_L$ is sufficiently small.

Let $p=\Pi_{\le r}m$. On the squared-coercivity event,
training consistency and the RKHS evaluation inequality imply
\[
 \frac14\|p\|_2^2
 \le\frac1n\sum_i(f(X_i)p(X_i))_-^2
 \le\frac1n\sum_i(m(X_i)-p(X_i))^2
 \le\frac{\|m\|_K^2}{n}\|Q_{r,SS}\|_{\mathrm{op}}
 \le b_r\tau_r.
\]
Choose an odd $u=r+O(\log d)$, with a sufficiently large fixed implicit
constant. Compare the sharp eigenvalue bounds at the two endpoints:
$\ell_u/\ell_r\le C_L\binom{d+1}{r+2}/\binom{d+1}{u+2}$,
up to a bounded power ratio. The binomial ratio is geometrically small
in $u-r$, uniformly for $r\le\alpha d$. The sharp level masses
and the bound $r\ge(\log d)^2/4$ therefore give
\begin{equation}\label{eq:negative-high-cutoff}
 \tau_u=(1-o(1))\tau_r,\qquad n\ell_u=o(\tau_r).
\end{equation}
These estimates are uniform over the sample-size range under
consideration: the removed band has trace $O((\log d)/r)\tau_r$,
and the geometric decay over $O(\log d)$ levels absorbs all polynomial
factors in \eqref{eq:negative-cutoff-bounds}.
Put $h=\Pi_{\le u}m$ and $t=m-h$. Orthogonality yields
\begin{equation}\label{eq:negative-low-norm}
 \|h\|_2^2\le4b_r\tau_r+n\ell_r
             \le C_{L,A}\tau_r(b_r+L_\rho),
 \qquad \|t\|_2^2=o(\tau_r).
\end{equation}

Lemma~\ref{lem:negative-mean-confidence} and
\eqref{eq:negative-high-cutoff} show that outside an $o(1)$ fraction
of inputs, $\sigma_x^2\ge\tau_r/3$ and
$|t(x)|\le c_\theta\sqrt{\tau_r}/(2\sqrt3)$.
Set $a_r=c_\theta\sqrt{\tau_r}/(2\sqrt3)$.
At every remaining input where $q_S(x)>\theta$, we have
$f(x)h(x)>a_r$. Consequently
\begin{equation}\label{eq:negative-confidence-clipping}
 \mathbb E f\,\operatorname{clip}(h/a_r,-1,1)
 \ge1-2F_\theta(S)-o(1).
\end{equation}
For any polynomial $h$ of degree at most $u$,
Lemma~\ref{lem:clipped-parity} gives
\[
 \left|\mathbb E f\,\operatorname{clip}(h/a_r,-1,1)\right|
 \le\frac{\|\Delta h\|_1}{a_r d}
 \le\frac{u\|h\|_2}{a_r d},
 \qquad \Delta=\sum_{j=1}^dD_j.
\]
For clarity, the first inequality follows directly by setting $q=fh$
and letting $R$ average the one-coordinate flips. Then
$Rq=-q+(2/d)f\Delta h$. Stationarity, Jensen's inequality and the
triangle inequality imply
$|\mathbb E|q+a_r|-\mathbb E|q-a_r||
\le(2/d)\|\Delta h\|_1$.
Use $\operatorname{clip}(q/a_r,-1,1)
=(|q+a_r|-|q-a_r|)/(2a_r)$.
The second inequality uses the Walsh multiplier bound
$\|\Delta h\|_2\le u\|h\|_2$.
Thus \eqref{eq:negative-low-norm} bounds the absolute value in
\eqref{eq:negative-confidence-clipping} by
\[
 C_{L,A,c_\theta}\rho
       \sqrt{1+\sqrt{L_\rho/\rho}+L_\rho}+o(1),
\]
which tends to zero as $\rho\downarrow0$. Choose $\alpha$ small enough
that it is at most $1/8$ for $0<\rho\le\alpha$, and then take $d$
large enough to absorb the vanishing errors. This proves
$F_\theta(S)\ge3/8$. The sampling and Gram events have combined
probability at least $3/4$ after the initial fixed choices.

It remains to cover smaller samples, including cases where the cutoff
$r$ does not exist. If $n<\exp((\log(d+1))^3)$, choose an odd
$v\asymp d^{1/3}$. The sharp spectrum gives
$n\ell_v/\tau_v=o(1)$, because
$\log\binom{d+1}{v+2}\asymp v\log(d/v)$ dominates $\log(n+1)$.
Let $p=\Pi_{\le v}m$. By
Lemma~\ref{lem:negative-mean-confidence}, outside an $o(1)$ fraction
of inputs, the residual variance is at least $\tau_v/2$ and
$|m-p|<c_\theta\sqrt{\tau_v}/(2\sqrt2)$.
Therefore confidence above $\theta$ implies $fp>0$ there.
Lemma~\ref{lem:ptf-parity}, with $v=o(\sqrt d)$, gives
$F_\theta(S)\ge1/2-o(1)\ge3/8$ for every such training sample.
For $n=0$ the conclusion also follows immediately from $q_S=1/2$.
Finally, the conditional width limit and $\theta<2/3$ give the asserted
finite-width $\liminf$ statement.
\end{proof}
